\section{Timing helpers}
\label{sec:timing}

The stateful timing helpers under \texttt{Helpers/Timing} are ordinary C\# objects. The caller supplies time, which keeps them useful in both Unity update loops and tests.

\subsection{Cooldown}
\texttt{Cooldown} exposes \texttt{Duration}, \texttt{Remaining}, \texttt{CanConsume}, and \texttt{IsActive}. Call \texttt{TryConsume()} to start the cooldown only when it is ready, then call \texttt{Update(deltaTime)} each frame.

\begin{lstlisting}
var attackCooldown = new Cooldown(0.5f);

if (attackCooldown.TryConsume())
{
    Attack();
}

attackCooldown.Update(Time.deltaTime);
\end{lstlisting}

\texttt{Start()} restarts the configured duration and \texttt{Reset()} returns the cooldown to its ready state.

\subsection{RateLimiter}
\texttt{RateLimiter} limits acquisitions in a rolling window. Construct it with a maximum number of uses and a window length, then pass a monotonic timestamp to \texttt{TryAcquire(double)}.

\begin{lstlisting}
var limiter = new RateLimiter(maxUses: 3, windowSeconds: 1.0);
double now = Time.realtimeSinceStartupAsDouble;

if (limiter.TryAcquire(now))
{
    SendRequest();
}
\end{lstlisting}

Timestamps must be finite and must not move backwards. \texttt{Reset()} clears the recorded acquisition times.
